\scanpage{140}

This, together with the relations
\[
H^0(\widehat{\mathbb Z}^{*}\times\mathfrak S_3,\widehat M)=H^1(\widehat{\mathbb Z}^{*}\times\mathfrak S_3,\widehat M)=0,
\]
made it possible to interpret the group $G_{\rho,\widetilde\sigma}$ as a group of transformations on the set $E$ of its own splittings, which made it possible to interpret the homomorphism $\mathcal M\to G_{\rho,\widetilde\sigma}$ (or $\widetilde{\mathcal M}\to G_{\rho,\widetilde\sigma}$) as an \deleted{\unsure{ho}} \uline{operation} of $\mathcal M$ \deleted{\unsure{of $H$}} (or better, \deleted{of $\unsure{\widetilde{\mathcal M}}$}\inserted{of $\widetilde{\mathcal M}=\pi_1(U_{0,3\mathbb Q},\mathfrak S_3)=\pi_1(M'_{1,1\mathbb Q})$} on $E$), thus corresponding to a \inserted{pro-étale} covering of $(U_{0,3\mathbb Q},\mathfrak S_3)\simeq M'_{1,1\mathbb Q}$ --- and this covering is none other than the Fermat pro-covering.

It remained to interpret in an analogous way the homomorphism $\boxed{\mathcal M=N_{1,1}\longrightarrow G_{\rho,\sigma}}$, fitting into the homomorphism of exact sequences
\[
\tag{298}
\diagram{\begin{tikzcd}[ampersand replacement=\&,column sep=small,row sep=large]
1\arrow[r]\&\begin{gathered}\mu_{\mathcal M}\\\Vert\\(\{1,\omega\})\end{gathered}\arrow[r]\arrow[d,"\simeq"]\&\begin{gathered}\mathcal M\\\Vert\\(N_{1,1})\end{gathered}\arrow[r]\arrow[d]\&\begin{gathered}\widetilde{\mathcal M}\\\Vert\\(N'_{1,1}=\unsure{\mathcal M_{0,3}})\end{gathered}\arrow[r]\arrow[d]\&1\\
1\arrow[r]\&\begin{gathered}\mu_G\\\Vert\\\{1,\omega\}\end{gathered}\arrow[r]\&G_{\rho,\sigma}\arrow[r]\&G_{\rho,\widetilde\sigma}\arrow[r]\&1
\end{tikzcd}},
\]
\deleted{For} which moreover makes it possible to interpret $\mathcal M$, as a central extension of $\widetilde{\mathcal M}$
\margin{139}

% Archive printed footer (not manuscript):
% 139
% https://grothendieck.umontpellier.fr
% Université de Montpellier, UMR CNRS 5149, Case courrier 051, Place Eugène Bataillon
% 34095 Montpellier cedex 5 - France

\scanpage{141}
\mspage{718}

by $\mu$, as the inverse image of the extension $G_{\rho,\sigma}$ of $G_{\rho,\widetilde\sigma}$ by $\mu$, under the preceding homomorphism $\widetilde{\mathcal M}\to G_{\rho,\widetilde\sigma}$. There is a striking analogy with what had happened, moreover, when we had started \unsure{from} (in \S~48, and \unsure{in|of} certain developments which preceded it) from the quotient group (``of Lie type'') $\unsure{\mathrm{SL}(2,\widehat{\mathbb Z})'}$ of $\unsure{\mathfrak S_{0,3}^{+\wedge}}$, to construct a homomorphism
\[
\deleted{\unsure{\mathcal M}}\qquad \mathcal M_{0,3}\longrightarrow\mathrm{GL}(2,\widehat{\mathbb Z})/\{\pm1\}
\]
from which it was possible for us to reconstruct \deleted{$\unsure{\mathcal M}$} the extension $\mathcal M=N_{1,1}$ of $\mathcal M_{0,3}$ by $\mu$ precisely as the inverse image of the extension $\mathrm{GL}(2,\widehat{\mathbb Z})$ of $\mathrm{GL}(2,\widehat{\mathbb Z})/\{\pm1\}$ by $\mu=\{\pm1\}$. From the present point of view where I am trying to make arithmetic-geometric aspects as explicit as possible in terms of automorphisms of the \inserted{profinite} loop groups $\Pi_{0,3}^\wedge$ --- the remarkable fact is that the two
\margin{140}

% Archive printed footer (not manuscript):
% 140
% https://grothendieck.umontpellier.fr
% Université de Montpellier, UMR CNRS 5149, Case courrier 051, Place Eugène Bataillon
% 34095 Montpellier cedex 5 - France

\scanpage{142}

extensions of $\widetilde{\mathcal M}$ by $\mu$, coming out of apparently rather different calculations, are canonically isomorphic.

I presume that a general explanation of this fact \unsure{will not present|would not present} any difficulty, in the present context --- because in both cases, the \inserted{bi-}\unsure{covering} of $\unsure{\mathrm{GL}|\mathrm{SL}}(2,\widehat{\mathbb Z})/\pm1$ resp. of $G_{\rho,\widetilde\sigma}$ came from the fact that I had started from a quotient group ($\mathrm{SL}(2,\widehat{\mathbb Z})$ or $\mathfrak S^{+\wedge}$) \uline{not} of $\mathfrak S_{0,3}^{+\wedge}\simeq\mathrm{SL}(2,\mathbb Z)^\wedge/\pm1$, but of $\mathrm{SL}(2,\mathbb Z)^\wedge\simeq N_{1,1}$ itself ---

If we nevertheless look for an interpretation (of the ``Fermat tower'' kind) of the homomorphism $\mathcal M\to G_{\rho,\sigma}$, it would be necessary to interpret $G_{\rho,\sigma}$ \inserted{(as \unsure{was} $G_{\rho,\widetilde\sigma}$)} as a group of automorphisms of a ``typical fiber'' ``not too big'' --- the difficulty being to invent something on which \unsure{that|it} operates nontrivially.
\margin{141}

% Archive printed footer (not manuscript):
% 141
% https://grothendieck.umontpellier.fr
% Université de Montpellier, UMR CNRS 5149, Case courrier 051, Place Eugène Bataillon
% 34095 Montpellier cedex 5 - France

\scanpage{143}
\mspage{719}

Interpreting $\mathcal M_{0,3}=N'_{1,1}$ as $\pi_1(M'_{1,1})$, where $M'_{1,1}$ is the modular scheme/$\mathbb Q$ of elliptic curves up to symmetry, we see that the Fermat tower amounts to this: over every base scheme $S$ (say \inserted{\uline{of characteristic 0}}) and \unsure{every|every} quasi-elliptic curve \inserted{\deleted{$\unsure{E}$}$\widetilde X$} over $S$, let us consider the corresponding étale covering of degree three $S'$ of $S$ (\inserted{of the} ``points of order 2 of $\widetilde X$''), \deleted{let us put} let
\[
\tag{299}
\deleted{\unsure{M_2}}\quad M_0(\widetilde X/\unsure{S|B})\overset{\mathrm{dfn}}{=}p_*(\mathbb Z_{\unsure{S'|B'}})/\mathbb Z_S
\]
be the corresponding twisted constant $\mathbb Z$-module free of rank \deleted{three}\inserted{two}, and let us consider
\[
\tag{300}
\deleted{\unsure{M_2}}\quad M_1(\widetilde X/S)\overset{\mathrm{dfn}}{=}M_0(\widetilde X/S)\otimes_{\mathbb Z}\underbrace{\unsure{\widehat{\mathbb Z}(1)}}_{\inserted{\begin{gathered}\Vert\,\mathrm{dfn}\\T_\infty(\mathbb G_m)\end{gathered}}},
\]
which is a pro-étale \deleted{pro}group over $S$, a system of coefficients locally isomorphic (in the broad sense) to $\widehat{\mathbb Z}^{2}$. \deleted{This being so,}\inserted{Of} course, $M_1(\widetilde X/S)$ depends \unsure{functorially} on $\widetilde X/S$ for isomorphisms and inverse images (we have a cartesian section of an ad-hoc fibered category $\ldots$). This being so,
\margin{142}

% Archive printed footer (not manuscript):
% 142
% https://grothendieck.umontpellier.fr
% Université de Montpellier, UMR CNRS 5149, Case courrier 051, Place Eugène Bataillon
% 34095 Montpellier cedex 5 - France

\scanpage{144}

giving the ``Fermat tower'' is equivalent to giving, for $\forall\,\widetilde X/S$ as above, a torsor
\[
\tag{301}
F_\infty(\widetilde X/S)\quad\text{torsor over }S\text{ under }M_1(\widetilde X/S),
\]
\margin{for a geometric construction cf. (320) page 724}
compatible with isomorphisms and \unsure{changes} of base.

The preceding theorem says essentially that such an $F_\infty(\widetilde X/S)$ exists, and is unique up to unique isomorphism, requiring in addition that the natural homomorphism
\diagram{\begin{tikzcd}[ampersand replacement=\&]
\Pi_{0,3}^{\wedge}\arrow[r] \& M^{\wedge}
\end{tikzcd}}
associated with such a datum $F_\infty$ be the canonical homomorphism $\ldots$.

On the other hand, let $X$ be a relative elliptic curve over a base $S$ of characteristic 0, and let $S'\subset X$ be the closed subscheme, étale of degree 3 over $S$, corresponding to the points of order 2. Let us consider the \unsure{generalized Jacobian} of $X$ deduced from this multisection, let $\mathcal J^0$ be its part of degree 0, we thus have

\margin{143}

% Archive printed footer (not manuscript):
% 143
% https://grothendieck.umontpellier.fr
% Université de Montpellier, UMR CNRS 5149, Case courrier 051, Place Eugène Bataillon
% 34095 Montpellier cedex 5 - France

\scanpage{145}
\mspage{720}

an exact sequence
\[
\tag{302}
\diagram{\begin{tikzcd}[ampersand replacement=\&]
1\arrow[r] \& \mathcal J_0^0\arrow[r,dash]\arrow[d,phantom,"\rotatebox{90}{$\simeq$}" description] \& \mathcal J^0\arrow[r] \& X\arrow[r,dash] \& 1\\
\& \mathbb G_m\otimes_{\mathbb Z}M_0
\end{tikzcd}}
\]
where $M_0=M_0(X/S)$ is defined above (299). From this we deduce an exact sequence
\[
\tag{303}
\diagram{\begin{tikzcd}[ampersand replacement=\&]
1\arrow[r] \& M_1\arrow[r]\arrow[d,equal] \& T_\infty(\mathcal J^0)\arrow[r] \& T_\infty(X)\arrow[r,dash] \& 1\\
\& M_0\otimes_{\mathbb Z}T_\infty(\mathbb G_m)
\end{tikzcd}}
\]
where
\[
\tag{\uline{304}}
T_\infty(G)=\varprojlim_n {}_nG
\]
is the local system of the $\Pi_1$ of the smooth commutative group scheme \inserted{with connected fibers} over $S$. Thus giving a torsor under $M_1$ is equivalent to giving a torsor under $T_{\unsure{\infty|\alpha}}(\mathcal J^0)$, equipped with a section of the associated torsor \unsure{of the|of} group $T_\infty(X)$.
\unsure{Is there|Is there} a remarkable description of the Fermat torsor $F_\infty(X/S)$, in terms of the geometry of the relative elliptic curve $X/S$?

\margin{144}

% Archive printed footer (not manuscript):
% 144
% https://grothendieck.umontpellier.fr
% Université de Montpellier, UMR CNRS 5149, Case courrier 051, Place Eugène Bataillon
% 34095 Montpellier cedex 5 - France

\scanpage{146}

$14^\circ)$ Relation to generalized Jacobians.

Let us consider the ``generalized Jacobian'' of
\[
\inserted{U_{0,3}={}}U_{0,3\,\deleted{\unsure{\mathbb Q}}\inserted{\mathbb Z}}
=\deleted{\unsure{\mathbb A_{\mathbb Q}}}\mathbb P^1_{\deleted{\unsure{\mathbb Q}}\inserted{\mathbb Z}}
-\{0,1,\infty\}_{\deleted{\unsure{\mathbb Q}}\inserted{\mathbb Z}},
\]
it is thus a group scheme $\mathcal J^*_{U_{0,3}}=\mathcal J^*$, \unsure{augmented|mapping} to $\mathbb Z$, fitting into the exact sequence
\[
\tag{305}
\diagram{\begin{tikzcd}[ampersand replacement=\&]
1\arrow[r] \& \mathcal J^0\arrow[r]\arrow[d,phantom,"\rotatebox{90}{$\simeq$}" description] \& \mathcal J^*\arrow[r,"\deg"] \& \mathbb Z_{\inserted{S_0}\deleted{\unsure{\mathbb Q}}}\arrow[r] \& 1\\
\& \mathbb G_{mS_0}\otimes_{\mathbb Z}M_{\unsure{0|\mathbb Z}}
\end{tikzcd}},
\]
\margin{\uline{N.B.} $S_0\overset{\unsure{\mathrm{df}|\mathrm{dfn}}}{=}\operatorname{Spec}\mathbb Z$}
and equipped with a canonical \deleted{homo}morphism
\[
\tag{306}
\diagram{\begin{tikzcd}[ampersand replacement=\&]
\varphi:\ U_{0,3\,\deleted{\unsure{\mathbb Q}}\inserted{\mathbb Z}}\arrow[r] \&
\underbrace{\mathcal J^1}_{\substack{\text{subscheme of }\mathcal J^*\\\text{inverse image of}\\\text{the section }1\text{ of }\mathbb Z_{S_0}}}
\end{tikzcd}}
\]
The operations of $\mathfrak S_3$ on $U_{0,3}$ define by functoriality operations of $\mathfrak S_3$ on $\mathcal J^*$, compatible with the extension structure (305), and such that (306) commutes with the operations of $\mathfrak S_3$.
Let us note that $\mathcal J^1$ is a $\mathcal J^0$-torsor and

\margin{145}

% Archive printed footer (not manuscript):
% 145
% https://grothendieck.umontpellier.fr
% Université de Montpellier, UMR CNRS 5149, Case courrier 051, Place Eugène Bataillon
% 34095 Montpellier cedex 5 - France